\documentclass[11pt]{article}

\usepackage{amsmath}
\usepackage{graphicx}
\usepackage{booktabs}
\usepackage{geometry}
\usepackage{hyperref}
\usepackage{multirow}
\usepackage{subcaption}

\geometry{margin=1in}

\title{CNN-Based Soil Moisture Forecasting Using TxSON In-Situ Data}
\author{Abijith Chandrasekar Vijayan}
\date{\today}

\begin{document}

\maketitle

% ============================================================================
\section{Introduction}
% ============================================================================

Soil moisture is a key variable in hydrology, agriculture, and climate science.
Accurate forecasts of near-surface soil water content help inform irrigation
scheduling, drought monitoring, and land--atmosphere interaction modeling. Because
soil moisture evolves over time---driven by precipitation, evapotranspiration, and
drainage---temporal forecasting methods are a natural fit for the problem: given a
recent history of observations, predict what soil moisture will be some number of
hours into the future.

Traditional approaches to soil-moisture forecasting include process-based land
surface models and statistical methods such as autoregressive models. More recently,
deep learning methods have been applied to this problem, motivated by their ability
to learn nonlinear temporal patterns directly from data without requiring
hand-specified physical equations. Among these, one-dimensional convolutional neural
networks (1D CNNs) have shown promise for time-series forecasting tasks
\cite{borovykh2017conditional, bai2018empirical}. A 1D CNN can learn short-term
patterns in the time series by applying the same learned filters across different
parts of the input window, which makes it well suited for detecting local events
like rain-driven moisture rises regardless of when they occur within the window.

This report describes CNN-based soil moisture forecasting using hourly in-situ
observations from the Texas Soil Observation Network (TxSON). The experiments
progress through three stages: (1)~a single-station baseline comparing two CNN
architectures against simple baselines across window sizes and horizons;
(2)~systematic feature engineering and architecture search, testing air temperature,
soil moisture at other depths, rainfall accumulations, and seasonal encoding;
(3)~multi-station generalization across four TxSON stations, evaluating both temporal
generalization (hold-out year) and spatial generalization (hold-out station).

In addition, this report investigates a classification formulation of the
forecasting problem: instead of predicting an exact soil moisture value, the CNN
predicts which of 12~fixed moisture categories the target falls into. Three
classification CNN architectures are compared against each other and against
non-learned baselines, using both standard and distance-aware evaluation metrics
that exploit the ordinal structure of the moisture categories.

% ============================================================================
\section{Dataset and Preprocessing}
% ============================================================================

\subsection{Station and Data Source}

All experiments use hourly observations from station \textbf{FD11} in the TxSON33
modeling-complete dataset. FD11 was selected because it has nearly complete hourly
coverage across a long time span, and the data are almost entirely observed-origin
for both the target variable and precipitation. A few years contain very small
amounts of \texttt{pipeline\_imputed} values, but no
\texttt{modeling\_completion} (synthetic) values are used for evaluation, so all
test-set targets are real sensor readings.

The dataset provides hourly measurements of soil water content at 20\,cm depth
(\texttt{SWC\_20}), precipitation (\texttt{Ppt}), and associated
\texttt{*\_Origin} columns that record whether each measurement is
\texttt{observed}, \texttt{pipeline\_imputed}, or \texttt{modeling\_completion}.
The dataset documentation explicitly notes that \texttt{modeling\_completion} values
are synthetic and should not normally be used as validation or test ground truth.

\subsection{Prediction Target}

The prediction target throughout is \texttt{SWC\_20}: soil water content at 20\,cm
depth, measured in $\text{m}^3/\text{m}^3$. This is consistent with the rest of the
project's existing forecasting work.

\subsection{Years Used}

Ten complete calendar years are available for FD11 (2015--2024; 2014 and 2025 are
partial years and are excluded). Three of these years are used as held-out test
years, chosen to span the wet/dry range by annual precipitation total:
\begin{itemize}
    \item \textbf{Driest:} 2023 (203\,mm annual precipitation)
    \item \textbf{Near-median:} 2020 (635\,mm; a leap year with 8{,}784 hours)
    \item \textbf{Wettest:} 2018 (856\,mm)
\end{itemize}
Validation is fixed at 2019 (mid-range precipitation, never used as a test year)
across all three folds. The remaining 8 years are pooled for training.

\subsection{Input Features}

For regression experiments, two feature configurations are compared:
\begin{enumerate}
    \item \textbf{SWC\_20 only} ($F = 1$): a univariate setting where the model sees
    only past soil moisture.
    \item \textbf{SWC\_20 + Ppt\_log} ($F = 2$): soil moisture plus log-transformed
    precipitation, $\texttt{Ppt\_log} = \log(1 + \texttt{Ppt})$. The log
    transformation compresses the long right tail of precipitation values (mostly
    zero, with occasional storms) into a more model-friendly range.
\end{enumerate}
Precipitation values are clipped at zero before the log transformation as a
defensive step against potential negative sensor noise, though no negative values
were present in the FD11 data used here.

For classification experiments, the feature sets are SWC\_20~only and
\textbf{SWC\_20 + Tair} ($F = 2$), where Tair is air temperature. This choice
follows from the feature engineering results (Section~5.2), which identified Tair
as the only feature to consistently improve CNN performance at longer horizons.

\subsection{Input Windows and Forecast Horizons}

Three input-window sizes are tested: 12, 24, and 48 hours. Each window is a
contiguous block of hourly observations that serves as the model's input.

Four forecast horizons are evaluated:
\begin{itemize}
    \item 24 hours (1 day ahead)
    \item 48 hours (2 days ahead)
    \item 72 hours (3 days ahead)
    \item 168 hours (1 week ahead)
\end{itemize}
Each horizon specifies how many hours after the window's last observed hour the
target value sits. Windows are built independently within each calendar year---a
window's input hours and its target hour must fall in the same year, so no window
ever spans a year boundary. This prevents any leakage across the held-out year
boundary by construction, at the cost of discarding a small number of windows near
each year's end.

\subsection{Normalization}

All input features are normalized to approximately $[0, 1]$ using a
\texttt{MinMaxScaler} fit on the pooled training years only. The same fitted scaler
is then applied (without refitting) to transform the validation and test data. This
means validation and test values can legitimately fall outside $[0, 1]$ if those
periods are wetter or drier than the training years---this is expected behavior, not
a bug.

Fitting the scaler only on training data matters because doing otherwise would leak
information about the validation or test distribution into the preprocessing step,
making results artificially well-behaved.

\subsection{Inverse Transformation for Metric Reporting}

All evaluation metrics (MSE, MAE, SMAPE, RSE) are computed on the original
\texttt{SWC\_20} scale, not the internal normalized scale. After training and
prediction (which occur entirely in scaled space), both true values and predictions
are inverse-transformed back to physical units using each fold's own scaler before
any metric is calculated. This is necessary because each fold's scaler has slightly
different min/max parameters (different training years span different moisture
ranges), so a scaled-space error of, say, 0.001 in one fold is not the same physical
error as 0.001 in another fold. Pearson correlation (CORR) is not affected by this
inverse transformation, since correlation does not change when values are linearly
rescaled.

% ============================================================================
\section{Experimental Setup}
% ============================================================================

\subsection{Year-Wise Held-Out Evaluation}

Each of the three test years is held out one at a time. For each fold, the model is
trained on the remaining 8 years (the training set), and validation year 2019 is used
solely for monitoring validation loss and making early-stopping decisions---it is not
part of the training data. The held-out test year is never seen during training or
validation and is used only for final metric computation.

\subsection{Baseline Models}

Two non-learned baselines are evaluated with the exact same metrics and test windows
as the CNNs:
\begin{itemize}
    \item \textbf{Persistence baseline:} predicts that soil moisture at the target
    hour equals the last observed value in the input window.
    \item \textbf{Moving-average baseline:} predicts the mean of \texttt{SWC\_20}
    across the entire input window. The window size used by the moving average
    matches the window size being tested in each configuration (12, 24, or 48
    hours), keeping the comparison fair.
\end{itemize}
Neither baseline has learned parameters. They are included specifically to determine
whether the CNN's added complexity provides any benefit beyond these simple
approaches.

\subsection{Model Comparison}

Two CNN architectures are compared:
\begin{itemize}
    \item A \textbf{small CNN} with 73 trainable parameters (one convolutional layer,
    8 filters).
    \item A \textbf{larger CNN} with 2{,}177 trainable parameters (two convolutional
    layers, 16 and 32 filters).
\end{itemize}
Both use the same training procedure, loss function, and evaluation pipeline, so any
difference in results is attributable to model capacity rather than training
configuration.

\subsection{Feature-Set Comparison}

Each CNN configuration is tested with both \texttt{SWC\_20} only and
\texttt{SWC\_20 + Ppt\_log} to determine whether adding precipitation as an input
improves forecasting performance.

The full experiment matrix consists of 3 test years $\times$ 3 window sizes $\times$
4 horizons $\times$ 2 feature sets $\times$ 4 models (2 baselines + 2 CNNs) = 288
rows, though only \textbf{216 unique experiments}: the two baselines do not use
precipitation, so their rows under each feature set are exact duplicates.

\subsection{Feature Engineering}

Six candidate features are tested individually against the SWC\_20-only baseline:
\begin{itemize}
    \item \textbf{Ppt\_log}: $\log(1 + \texttt{Ppt})$, log-transformed precipitation.
    \item \textbf{Tair}: air temperature.
    \item \textbf{T\_20}: soil temperature at 20\,cm depth.
    \item \textbf{SWC depths}: soil moisture at 5, 10, and 50\,cm
    (\texttt{SWC\_5}, \texttt{SWC\_10}, \texttt{SWC\_50}).
    \item \textbf{Rain accumulations}: backward-looking 3h, 6h, 12h, and 24h rolling
    sums of precipitation.
    \item \textbf{Seasonal encoding}: $\sin(2\pi \cdot \text{doy}/365.25)$ and
    $\cos(2\pi \cdot \text{doy}/365.25)$, providing a smooth cyclical representation
    of the time of year where December~31 is close to January~1.
\end{itemize}
Selection is done on \textbf{validation loss} (2019 data), not test performance, to
avoid data leakage into the feature selection process.

\subsection{Multi-Station Generalization}

Four TxSON stations sharing 10~complete calendar years (2015--2024) are used for
generalization experiments: FD11, FD02, CB06, and CB25. Two experiment types are
conducted:
\begin{itemize}
    \item \textbf{Experiment~A (hold-out year):} All 4 stations contribute training
    data; one calendar year is held out from all stations for testing.
    \item \textbf{Experiment~B (hold-out station):} One station is entirely removed
    from training; the model trains on the other 3 stations and is tested on all
    non-validation years of the held-out station.
\end{itemize}
In both cases, validation is fixed at 2019 across all stations. The scaler is fit on
pooled training data only, and windows are built independently within each
(station, year) pair.

\subsection{12-Class Classification}

As a separate line of investigation, the continuous SWC\_20 target is discretized
into 12~fixed categories of width 0.05:
\begin{equation}
    \text{Class } c = \lfloor \texttt{SWC\_20} / 0.05 \rfloor, \quad c \in \{0, 1, \ldots, 11\}
\end{equation}
covering the range $[0.00, 0.60]$. The same bin boundaries are applied to every
station, so that class~5, for example, always means
$\text{SWC\_20} \in [0.25, 0.30)$. A classification CNN replaces the single
regression output with 12~logits, trained with cross-entropy loss. The same input
pipeline, window construction, and train/val/test splits are reused.

Three classification CNN architectures are compared:
\begin{itemize}
    \item \textbf{SmallClassCNN:} 1~convolutional layer, 8~filters, average pooling
    (128~params with $F=1$).
    \item \textbf{LargerClassCNN:} 2~convolutional layers, 16/32 filters, average
    pooling (2{,}364~params with $F=1$).
    \item \textbf{MaxpoolClassCNN:} Same as LargerClassCNN but with max pooling
    instead of average pooling (2{,}364~params with $F=1$).
\end{itemize}
All three share the same fully connected head: pooled features $\to$ 16-unit hidden
layer with ReLU $\to$ 12-class logit output. Two feature sets are tested:
SWC\_20~only ($F=1$) and SWC\_20 + Tair ($F=2$), at horizons H=24 and H=168.

\textbf{Loss functions.} Two loss functions are compared: standard cross-entropy and
class-weighted cross-entropy. For the weighted variant, inverse-frequency weights are
computed from the \emph{training} data only:
\begin{equation}
    w_c = \frac{N_{\text{train}}}{K_{\text{active}} \cdot n_c}
\end{equation}
where $n_c$ is the training count for class~$c$, $K_{\text{active}}$ is the number
of non-empty classes, and the weights are normalized so their mean equals~1. Empty
classes receive weight~0. This upweights rare classes (e.g., near the tails of the
SWC\_20 distribution) without introducing signal from validation or test data.

\textbf{Baselines.} Two non-learned classification baselines are evaluated:
\begin{itemize}
    \item \textbf{Persistence:} assigns the class of the last observed SWC\_20 value
    in the input window.
    \item \textbf{Majority class:} always predicts the most frequent class in the
    training data.
\end{itemize}

\textbf{Classification metrics.} Standard classification metrics (accuracy, macro
and weighted F1, per-class precision/recall) are supplemented with
\emph{distance-aware} metrics that exploit the ordinal structure of the 12~classes:
\begin{itemize}
    \item \textbf{Within-1-category accuracy:} fraction of predictions within one
    class of the true class.
    \item \textbf{Within-2-category accuracy:} fraction within two classes.
    \item \textbf{Mean absolute class error (MACE):} mean of
    $|c_{\text{pred}} - c_{\text{true}}|$ across test samples.
\end{itemize}
These metrics reflect the fact that misclassifying by one bin ($\pm 0.05$ SWC) is
far less consequential than misclassifying by three or more bins.

% ============================================================================
\section{CNN Architecture}
% ============================================================================

This section describes both CNN architectures mathematically. Notation follows the
general convention used for 1D convolutional sequence models in the time-series
forecasting literature \cite{borovykh2017conditional, bai2018empirical}, adapted to
this specific model.

\subsection{Input Representation}

Let $W$ be the input-window length in hours and $F$ the number of input features per
hour. A single input example is a matrix
\begin{equation}
    X \in \mathbb{R}^{W \times F}, \qquad
    X = (x_1, x_2, \dots, x_W), \quad x_t \in \mathbb{R}^F
\end{equation}
where $x_t$ is the feature vector observed at hour $t$ within the window. For the
\texttt{SWC\_20}-only setting, $F = 1$ and $x_t$ is a scalar; with precipitation
added, $F = 2$. PyTorch's \texttt{Conv1d} expects channels before the time axis, so
$X$ is transposed to $X^\top \in \mathbb{R}^{F \times W}$ before the first
convolution---this is just a reshaping step, not a computation.

\subsection{1D Convolution}

A convolution layer with $C_{\text{out}}$ output channels and kernel size $K$
produces, for output channel $c$ and time position $t$:
\begin{equation}
    z_{c,t} = \sum_{f=1}^{F} \sum_{k=1}^{K}
    w_{c,f,k} \; x_{f,\, t + k - \lceil K/2 \rceil} + b_c
\end{equation}
where $w_{c,f,k}$ are the learned filter weights and $b_c$ the learned bias for
output channel $c$. Each output channel is a learned pattern detector: a small
template of length $K$ across all $F$ input features. Its response $z_{c,t}$ is
large wherever the input around hour $t$ matches that template.

The key property of convolution is that the same filter is applied at every position
in the window. This means a pattern like a rain-driven moisture rise is detected the
same way whether it happens near the start or end of the window, rather than
requiring the network to learn it separately at every position.

\textbf{Padding.} With padding $= \lfloor K/2 \rfloor$, the input is zero-padded on
both sides so the output length stays $W$ rather than shrinking to $W - K + 1$. This
keeps the downstream pooling step operating over the full window regardless of $K$.

\subsection{Activation Function}

Each $z_{c,t}$ passes through the ReLU activation:
\begin{equation}
    a_{c,t} = \max(0,\; z_{c,t})
\end{equation}
applied elementwise. ReLU introduces nonlinearity---without it, stacking convolutions
and linear layers would collapse into a single linear function of the input, no
matter how many layers are used.

\subsection{Second Convolutional Layer (Larger CNN Only)}

The small model uses a single convolution ($F \to 8$ channels). The larger model
stacks two convolutions: $F \to 16 \to 32$ channels. The second convolution does
not see the raw input---it operates on the 16 filter responses produced by the first
layer. This lets it detect slightly higher-level patterns built on top of the first
layer's detections (for example, two short moisture changes happening close together
in time), which a single convolution layer cannot represent directly.

\subsection{Average Pooling}

\begin{equation}
    p_c = \frac{1}{W} \sum_{t=1}^{W} a_{c,t}
\end{equation}

Average pooling collapses the time axis entirely, producing one number per channel:
how strongly, on average across the whole window, that channel's pattern was present.

This is also where positional information is lost. The pooled value $p_c$ cannot
tell the difference between a pattern that appeared once strongly near the end of the
window and one that appeared at moderate strength throughout. This is a deliberate
simplification that keeps the model small, at the cost of not being able to represent
\emph{when} within the window something happened.

\subsection{Fully Connected Layers}

The per-channel pooled values are collected into a vector
$p \in \mathbb{R}^{C_{\text{out}}}$ (8-dimensional for the small model,
32-dimensional for the larger one). Two linear layers map this to a prediction:
\begin{align}
    h &= \text{ReLU}(W_1 \, p + b_1) \\
    \hat{y} &= W_2 \, h + b_2
\end{align}
where $W_1 \in \mathbb{R}^{H \times C_{\text{out}}}$ and
$W_2 \in \mathbb{R}^{1 \times H}$, with $H = 4$ for the small model and $H = 16$
for the larger one.

Two linear layers rather than a single $C_{\text{out}} \to 1$ mapping are used
because the relationship between ``how much of each pattern fired'' and the final
soil-moisture prediction is generally not linear. For example, a rising-moisture
signal combined with a low-baseline signal might matter differently together than
either one alone, and a single linear layer cannot capture that kind of interaction.

The final output $\hat{y} \in \mathbb{R}$ is the predicted (scaled) \texttt{SWC\_20}
value at $t_{\text{window\_end}} + \text{horizon}$.

\subsection{Architecture Summary}

Table~\ref{tab:arch} summarizes both architectures. The parameter counts assume $F=1$
(SWC\_20 only); with $F=2$, the first convolutional layer gains additional weights.

\begin{table}[h]
\centering
\caption{CNN architecture summary. Regression models output a single scalar;
classification models output 12~logits. LargerClassCNN (not shown) is identical
to MaxpoolClassCNN except it uses average pooling (2{,}364~params at $F\!=\!1$).}
\label{tab:arch}
\begin{tabular}{lllll}
\toprule
\textbf{Layer} & \textbf{SmallCNN} & \textbf{LargerCNN}
    & \textbf{SmallClassCNN} & \textbf{MaxpoolClassCNN} \\
\midrule
Conv1d \#1 & $F \to 8$, $K=3$ & $F \to 16$, $K=3$
    & $F \to 8$, $K=3$ & $F \to 16$, $K=3$ \\
Conv1d \#2 & --- & $16 \to 32$, $K=3$
    & --- & $16 \to 32$, $K=3$ \\
Pooling & AvgPool & AvgPool & AvgPool & \textbf{MaxPool} \\
FC \#1 & $8 \to 4$ + ReLU & $32 \to 16$ + ReLU
    & $8 \to 16$ + ReLU & $32 \to 16$ + ReLU \\
FC \#2 & $4 \to 1$ & $16 \to 1$
    & $16 \to 12$ & $16 \to 12$ \\
\midrule
Params ($F\!=\!1$) & 73 & 2{,}177 & 128 & 2{,}364 \\
\bottomrule
\end{tabular}
\end{table}

\subsection{Training Objective}

Regression models are trained by minimizing mean squared error (MSE):
\begin{equation}
    \mathcal{L}_{\text{reg}}(\theta) = \frac{1}{N} \sum_{i=1}^{N}
    \left( \hat{y}_i(\theta) - y_i \right)^2
\end{equation}
Classification models minimize cross-entropy loss over the 12~classes:
\begin{equation}
    \mathcal{L}_{\text{cls}}(\theta) = -\frac{1}{N} \sum_{i=1}^{N}
    w_{c_i} \log \hat{p}_{i,c_i}(\theta)
\end{equation}
where $c_i$ is the true class, $\hat{p}_{i,c_i}$ is the predicted probability for
that class, and $w_{c_i}$ is the class weight (all~1 for standard CE; inverse-frequency
weights for weighted CE). Classification models use 40~epochs with patience~8, compared
to 30~epochs with patience~6 for regression.

\subsection{Evaluation Metrics}

For $N$ test examples with true values $y_i$ and predictions $\hat{y}_i$:

\begin{align}
    \text{MSE} &= \frac{1}{N} \sum_{i} (\hat{y}_i - y_i)^2 \\[4pt]
    \text{MAE} &= \frac{1}{N} \sum_{i} |\hat{y}_i - y_i| \\[4pt]
    \text{SMAPE} &= \frac{100}{N} \sum_{i}
    \frac{2 |\hat{y}_i - y_i|}{|y_i| + |\hat{y}_i| + \varepsilon} \\[4pt]
    \text{RSE} &= \frac{\sum_{i} (\hat{y}_i - y_i)^2}
    {\sum_{i} (y_i - \bar{y})^2 + \varepsilon} \\[4pt]
    \text{CORR} &= \frac{\sum_{i} (y_i - \bar{y})(\hat{y}_i - \bar{\hat{y}})}
    {\sqrt{\sum_{i} (y_i - \bar{y})^2}
    \;\sqrt{\sum_{i} (\hat{y}_i - \bar{\hat{y}})^2}}
\end{align}
with $\varepsilon = 10^{-8}$ added only to prevent division by zero. RSE normalizes
MSE by the target's own variance (a model predicting the mean scores
$\text{RSE} \approx 1$). CORR measures whether predictions track the true trend
regardless of scale. All metrics except CORR are computed on original-scale
\texttt{SWC\_20} values (see Section~2.7); CORR is not affected by the inverse
transformation since it is invariant to linear rescaling.

% ============================================================================
\section{Training Setup}
% ============================================================================

Every CNN run uses the same fixed training configuration, so that differences in
results come from the variable being tested (model size, window size, feature set),
not from inconsistent training:

\begin{itemize}
    \item \textbf{Optimizer:} Adam \cite{kingma2015adam} with learning rate $10^{-3}$.
    \item \textbf{Loss function:} Mean squared error against the true (scaled)
    \texttt{SWC\_20} value.
    \item \textbf{Batch size:} 128.
    \item \textbf{Epoch budget:} Up to 30 epochs.
    \item \textbf{Early stopping:} Patience of 6---training stops if validation loss
    has not improved for 6 consecutive epochs, and the model reverts to the weights
    from its best validation epoch. This keeps the number of epochs actually used
    honest: a run that converged in 9 epochs and one that needed all 30 are both
    reported as such.
    \item \textbf{Validation monitoring:} Validation loss is computed at the end of
    each epoch on the fixed validation year (2019) without gradient updates.
    \item \textbf{Best-model restoration:} After training ends (either by early
    stopping or by reaching the epoch budget), the model is restored to the weights
    from the epoch with the lowest validation loss.
    \item \textbf{Reproducibility:} Python's \texttt{random}, NumPy, and PyTorch (CPU
    and CUDA) are all seeded with a fixed seed before every run.
\end{itemize}
Classification experiments use the same configuration except for a longer epoch
budget (40~epochs) and correspondingly larger patience (8~epochs), since the
12-class output space requires more training iterations to converge.

A separate extended-training experiment was also conducted: 100 epochs with early
stopping disabled, on one representative test year (2020), both models, at the 24h
and 168h horizons. This was used to inspect whether and when overfitting occurs---see
the Overfitting Analysis subsection in the Discussion.

% ============================================================================
\section{Results}
% ============================================================================

\subsection{Single-Station Baseline Experiment}

The baseline experiment matrix consists of 3~test years $\times$ 3~window sizes
$\times$ 4~horizons $\times$ 2~feature sets $\times$ 4~models (2~baselines +
2~CNNs) = 288~stored rows, corresponding to 216~unique experiments (baseline rows are
duplicated across feature sets since persistence and moving average ignore
precipitation).

Table~\ref{tab:baseline_horizon} summarizes the central comparison: CNN performance
versus the two baselines across forecast horizons, using the reference configuration
of window size~24 and SWC\_20~only features, averaged over the three test years.

\begin{table}[h]
\centering
\caption{Test MSE by model and horizon (W=24, SWC\_20 only, averaged over 3 test
years). Lower is better.}
\label{tab:baseline_horizon}
\begin{tabular}{lcccc}
\toprule
\textbf{Model} & \textbf{H=24} & \textbf{H=48} & \textbf{H=72} & \textbf{H=168} \\
\midrule
Persistence        & 0.00038 & 0.00097 & 0.00143 & 0.00216 \\
Moving Average     & 0.00070 & 0.00107 & 0.00149 & 0.00217 \\
SmallCNN           & 0.00045 & 0.00097 & 0.00146 & 0.00200 \\
LargerCNN          & 0.00037 & 0.00097 & 0.00143 & 0.00202 \\
\bottomrule
\end{tabular}
\end{table}

At H=24, persistence and LargerCNN are essentially tied (0.00038 vs 0.00037). At
H=48 and H=72, persistence is tied with or marginally ahead of both CNNs. At H=168,
both CNNs clearly beat persistence---SmallCNN (0.00200) and LargerCNN (0.00202) vs
persistence's 0.00216.

\textbf{Window size.} For SmallCNN, the 12h window produces the lowest MSE at every
horizon. For LargerCNN, the 24h window wins at H=24 and H=48, while 12h wins at
H=168. A plausible interpretation is that the larger model has enough capacity to
extract useful signal from the extra context a 24h window provides at shorter
horizons, while the smaller model does better with less context to sift through.

\textbf{Precipitation.} Adding log-transformed precipitation (Ppt\_log) reduced test
MSE in 6 of 8 model/horizon combinations: all 4 for SmallCNN, and 2 of 4 for
LargerCNN.

\textbf{Test-year wetness.} Among the three test years, 2018 (wettest, 856\,mm) was
consistently harder to predict than 2020 (median, 635\,mm) or 2023 (driest,
203\,mm), with roughly double the MSE across all horizons and both CNN architectures.
CORR was also consistently lower for 2018, indicating this is not simply an artifact
of larger variance in wet conditions.

\subsection{Feature Engineering}

Six candidate features were tested individually against the SWC\_20-only baseline at
H=24 (W=24, test year~2023, validation year~2019): log-precipitation, air temperature
(Tair), soil temperature at 20\,cm (T\_20), soil moisture at other depths
(SWC\_5/10/50), rolling rainfall accumulations (3h/6h/12h/24h), and seasonal encoding
(sin/cos of day-of-year).

No individual feature improved LargerCNN validation loss at H=24. This is consistent
with persistence being a strong baseline at short horizons: when the most recent
observation already contains most of the predictive signal for one day ahead,
additional features provide diminishing returns.

At longer horizons, air temperature (Tair) emerged as the only consistently useful
feature addition. Table~\ref{tab:feature_horizon} shows the validation loss
improvement from adding Tair at each horizon.

\begin{table}[h]
\centering
\caption{LargerCNN validation loss improvement from adding Tair (W=24, test=2023).}
\label{tab:feature_horizon}
\begin{tabular}{lcc}
\toprule
\textbf{Horizon} & \textbf{Val loss change} & \textbf{Direction} \\
\midrule
H=24  & $+0.1\%$  & No improvement \\
H=48  & $-3.9\%$  & Improved \\
H=72  & $-1.3\%$  & Improved \\
H=168 & $-8.0\%$  & Improved \\
\bottomrule
\end{tabular}
\end{table}

The trend is clear: Tair's value increases with the forecast horizon, consistent with
temperature being a driver of evapotranspiration processes that matter more over
multi-day timescales.

\textbf{Seasonal features.} Sine and cosine encoding of the day of year
($\sin(2\pi \cdot \text{doy}/365.25)$ and $\cos(2\pi \cdot \text{doy}/365.25)$) did
not improve validation loss at any horizon tested. This suggests that the CNN already
implicitly captures seasonal variation through the soil moisture values themselves.

\textbf{Combined feature set.} Combining multiple features (SWC at other depths,
Tair) caused training instability for LargerCNN at H=24, with a distribution-shift
diagnostic revealing that SWC\_5 and SWC\_10 have substantially different
distributions across the train/val/test splits.

\subsection{Architecture Search}

Ten CNN architectures were compared at H=24 and H=168, varying kernel size (3, 5, 7),
filter counts (8--64), depth (1--3 convolutional layers), pooling type (average vs
max), and dropout (0, 0.2, 0.3).

Key findings:
\begin{itemize}
    \item \textbf{Max pooling beats average pooling} at both H=24 and H=168, making
    it the single clearest architectural improvement identified.
    \item \textbf{Larger filter counts} (32/64) did not consistently improve test MSE
    over the baseline 16/32 configuration.
    \item \textbf{Dropout hurts at H=24}: adding dropout of 0.2 or 0.3 increased both
    validation loss and test MSE, likely because the models are already small enough
    that regularization is unnecessary.
    \item \textbf{Three convolutional layers} did not improve over two layers.
\end{itemize}

\subsection{Multi-Station Generalization}

The pipeline was extended to four TxSON stations that share 10~complete calendar years
(2015--2024): FD11, FD02, CB06, and CB25. Two experiment types test different kinds
of generalization:

\textbf{Experiment~A (hold-out year):} All 4 stations contribute training data, but
one calendar year is held out from all stations for testing. This tests temporal
generalization.

\textbf{Experiment~B (hold-out station):} One station is entirely removed from
training. The model trains on the other 3 stations and is tested on the held-out
station. This tests spatial generalization.

Both experiments use the same validation year (2019), window size (24), and two
horizons (H=24 and H=168). Models tested include persistence, moving average,
LargerCNN (average pooling), and MaxpoolCNN.

In Experiment~A, CNNs beat persistence at H=168 for 2020 and 2018 but not for 2023
(the driest year). At H=24, persistence remains competitive. Adding Tair improved CNN
performance at H=168, consistent with the single-station finding.

Experiment~B tests a harder form of generalization: predicting at a station the model
has never seen. Results show that spatial generalization is generally harder than
temporal generalization, with larger MSE increases when a station is held out compared
to when a year is held out.

\subsection{Class Distribution}

Before training, the distribution of the 12~classes was examined for each station.
At FD11, SWC\_20 ranges from 0.143 to 0.510, so classes~0, 1, and~11 are empty
(no samples). Classes~4--8 account for approximately 92\% of all samples, with
class~7 ($[0.35, 0.40)$) being the most frequent at 25.4\%. This imbalance
motivates the weighted cross-entropy investigation.

The four stations show substantially different moisture distributions under the
same fixed bin boundaries:
\begin{itemize}
    \item \textbf{FD11:} range $[0.143, 0.510]$, dominant class~7, empty
    classes~$\{0,1,11\}$.
    \item \textbf{FD02:} range $[0.064, 0.311]$, dominant class~2 (43.1\%), empty
    classes~$\{0,7,8,9,10,11\}$---a markedly drier station.
    \item \textbf{CB06:} range $[0.129, 0.457]$, dominant class~3 (32.0\%), empty
    classes~$\{0,1,10,11\}$.
    \item \textbf{CB25:} range $[0.096, 0.490]$, dominant class~4 (25.6\%), empty
    classes~$\{0,10,11\}$.
\end{itemize}
These differences mean that a multi-station classifier trained with fixed bins
must handle strongly heterogeneous class priors across stations.

\begin{figure}[h]
\centering
\includegraphics[width=\textwidth]{class_distribution_comparison.pdf}
\caption{Class distribution across four TxSON stations under the fixed 12-bin
scheme. FD02 is markedly drier (dominant class~2), while FD11 and CB25 span a
wider moisture range. Empty classes arise where a station's SWC\_20 values never
fall within a bin.}
\label{fig:cls_dist}
\end{figure}

\subsection{Single-Station Classification (FD11)}

Table~\ref{tab:cls_arch} compares the three classification architectures at H=24
and H=168, using SWC\_20-only inputs, averaged over the three test years.

\begin{table}[h]
\centering
\caption{Classification architecture comparison on FD11 (SWC\_20 only, W=24,
regular cross-entropy, averaged over 3 test years).}
\label{tab:cls_arch}
\begin{tabular}{llccccc}
\toprule
\textbf{H} & \textbf{Model} & \textbf{Acc} & \textbf{Macro F1}
    & \textbf{Wtd F1} & \textbf{W-1 Acc} & \textbf{MACE} \\
\midrule
24  & SmallClassCNN    & 0.866 & 0.663 & 0.860 & 0.979 & 0.169 \\
    & LargerClassCNN   & 0.871 & 0.669 & 0.866 & 0.980 & 0.162 \\
    & MaxpoolClassCNN  & 0.890 & 0.690 & 0.886 & 0.982 & 0.140 \\
    & Persistence      & 0.888 & 0.695 & 0.888 & 0.984 & 0.139 \\
    & Majority class   & 0.219 & 0.051 & 0.079 & 0.590 & 1.490 \\
\midrule
168 & SmallClassCNN    & 0.595 & 0.439 & 0.577 & 0.895 & 0.590 \\
    & LargerClassCNN   & 0.650 & 0.476 & 0.634 & 0.899 & 0.530 \\
    & MaxpoolClassCNN  & 0.651 & 0.477 & 0.634 & 0.901 & 0.525 \\
    & Persistence      & 0.584 & 0.442 & 0.584 & 0.917 & 0.560 \\
    & Majority class   & 0.217 & 0.051 & 0.078 & 0.589 & 1.498 \\
\bottomrule
\end{tabular}
\end{table}

At H=24, persistence is a strong baseline (accuracy 0.888, macro~F1 0.695),
consistent with the regression finding that the most recent observation is highly
predictive one day ahead. Among the three architectures, MaxpoolClassCNN achieves
the highest accuracy (0.890) and macro~F1 (0.690) at H=24, with LargerClassCNN
close behind (0.871, 0.669). The advantage of max pooling over average pooling
matches the regression-side finding that max pooling captures event-driven signals
more effectively.

At H=168, persistence degrades to 0.584 accuracy, and all three CNN architectures
clearly outperform it: MaxpoolClassCNN reaches 0.651 accuracy and LargerClassCNN
0.650, while even SmallClassCNN (0.595) exceeds persistence. The CNN--persistence
gap widens from near-zero at H=24 to $\sim$0.07 at H=168, paralleling the
regression experiments.

Within-1-category accuracy exceeds 0.979 at H=24 and remains above 0.895 at
H=168 for all architectures, indicating that even when the classifier errs, it is
almost always off by at most one bin (0.05~SWC units). Mean absolute class error
(MACE) ranges from 0.140 (MaxpoolCNN, H=24) to 0.590 (SmallCNN, H=168).

\begin{figure}[h]
\centering
\includegraphics[width=\textwidth]{cls_accuracy_by_arch.pdf}
\caption{Classification accuracy by architecture and horizon, averaged over three
test years (SWC\_20 only, regular CE). MaxpoolClassCNN is consistently best. The
persistence baseline is shown as a dashed line.}
\label{fig:cls_acc}
\end{figure}

\begin{figure}[h]
\centering
\includegraphics[width=\textwidth]{cls_confusion_matrices.pdf}
\caption{Confusion matrices for MaxpoolClassCNN on test year 2023 at H=24 (left)
and H=168 (right). At H=24, errors cluster tightly around the diagonal. At H=168,
the distribution spreads but remains within 1--2 bins of the true class.}
\label{fig:cls_cm}
\end{figure}

\subsection{Effect of Air Temperature}

Adding Tair as a second input feature ($F=2$) produces a mixed effect that varies
by horizon and test year. For MaxpoolClassCNN averaged over three test years,
SWC\_20+Tair slightly improves H=24 accuracy (0.895 vs.\ 0.890) and H=168 accuracy
(0.662 vs.\ 0.651). However, the benefit is year-dependent: Tair helps substantially
for 2018 (the wettest year; H=24 accuracy improves from 0.857 to 0.888) but hurts
for 2020 at H=24 (0.886 to 0.886, essentially flat) and at H=168 (0.729 to 0.683).
This suggests that temperature is most informative when evapotranspiration-driven
drying is an active process, as in wet years.

\subsection{Regular vs.\ Weighted Cross-Entropy}

Class-weighted cross-entropy upweights rare classes using inverse-frequency weights
computed from training data. Contrary to the expected tradeoff (improved macro~F1
at the cost of lower accuracy), weighted CE consistently degrades \emph{both}
accuracy and macro~F1 relative to standard CE across all three architectures and
both horizons.

The likely explanation is the extreme weight ratio in FD11's class distribution:
class~10, with only 6~training samples, receives a weight of 8.4, while the
dominant classes receive weights below 0.01. This 1{,}000$\times$ ratio
destabilizes gradient updates, causing the model to overfit to near-empty tail
classes at the expense of the well-populated middle classes that dominate the test
set. The effect is architecture-dependent: SmallClassCNN is most vulnerable
(3-year average accuracy drops from 0.595 to 0.465 at H=168), while MaxpoolClassCNN
is most robust (0.651 to 0.644, a negligible decline).

This result suggests that for distributions with extremely rare classes, naive
inverse-frequency weighting is counterproductive. More sophisticated approaches---such
as capping the maximum weight, using $\sqrt{1/n_c}$ instead of $1/n_c$, or focal
loss---may be needed to benefit from class reweighting without destabilizing training.

\subsection{Multi-Station Classification}

The classification framework was extended to the same four-station setting used for
the regression multi-station experiments (FD11, FD02, CB06, CB25). Two experiment
types mirror the regression design:

\textbf{Experiment~A (Hold-Out Year).} All four stations contribute data; one
calendar year is held out for testing, the preceding year for validation, and the
remaining years for training. Three test years (2018, 2020, 2023) are averaged.
Training sets contain $\approx$279{,}000 samples---roughly four times the
single-station equivalent.

\textbf{Experiment~B (Hold-Out Station).} One station is held out entirely for
testing (all its years), one year from the remaining three stations is used for
validation, and the rest for training. Each of the four stations takes a turn as
the hold-out.

Table~\ref{tab:multi_cls_a} reports Experiment~A results (SWC\_20 only, three-year
average). Multi-station training substantially improves temporal generalization
over single-station classification at both horizons. MaxpoolCNN reaches
0.927 accuracy at H=24 (vs.\ 0.890 single-station, +3.7~pp) and 0.708 at
H=168 (vs.\ 0.651, +5.7~pp). The larger and more diverse training set allows
the model to learn patterns that transfer across years more robustly.

\begin{table}[h]
\centering
\caption{Multi-station classification, Experiment~A: hold-out year
(SWC\_20 only, 3-year average). Best CNN values in bold.}
\label{tab:multi_cls_a}
\small
\begin{tabular}{llccccc}
\toprule
Horizon & Model & Acc & Macro F1 & Wtd F1 & W-1 Acc & MACE \\
\midrule
\multirow{5}{*}{H=24}
 & SmallCNN     & 0.899 & 0.713 & --- & 0.985 & --- \\
 & LargerCNN    & 0.898 & 0.717 & --- & 0.986 & --- \\
 & MaxpoolCNN   & \textbf{0.927} & \textbf{0.749} & --- & \textbf{0.988} & --- \\
 & Persistence  & --- & --- & --- & --- & --- \\
 & Majority     & --- & --- & --- & --- & --- \\
\midrule
\multirow{5}{*}{H=168}
 & SmallCNN     & 0.673 & 0.470 & --- & 0.928 & --- \\
 & LargerCNN    & 0.695 & 0.487 & --- & 0.929 & --- \\
 & MaxpoolCNN   & \textbf{0.708} & \textbf{0.506} & --- & \textbf{0.932} & --- \\
 & Persistence  & --- & --- & --- & --- & --- \\
 & Majority     & --- & --- & --- & --- & --- \\
\bottomrule
\end{tabular}
\end{table}

Table~\ref{tab:multi_cls_b} reports Experiment~B results (SWC\_20 only),
averaged over the four hold-out stations. Spatial generalization is harder than
temporal: MaxpoolCNN averages 0.885 accuracy at H=24 across the four stations
(vs.\ 0.927 in Experiment~A) and the gap widens at H=168.

\begin{table}[h]
\centering
\caption{Multi-station classification, Experiment~B: hold-out station
(SWC\_20 only, 4-station average). Best CNN values in bold.}
\label{tab:multi_cls_b}
\small
\begin{tabular}{llccccc}
\toprule
Horizon & Model & Acc & Macro F1 & Wtd F1 & W-1 Acc & MACE \\
\midrule
\multirow{5}{*}{H=24}
 & SmallCNN     & 0.840 & 0.608 & --- & 0.981 & --- \\
 & LargerCNN    & 0.855 & 0.594 & --- & 0.983 & --- \\
 & MaxpoolCNN   & \textbf{0.885} & \textbf{0.655} & --- & \textbf{0.985} & --- \\
 & Persistence  & --- & --- & --- & --- & --- \\
 & Majority     & --- & --- & --- & --- & --- \\
\midrule
\multirow{5}{*}{H=168}
 & SmallCNN     & --- & --- & --- & --- & --- \\
 & LargerCNN    & --- & --- & --- & --- & --- \\
 & MaxpoolCNN   & --- & --- & --- & --- & --- \\
 & Persistence  & --- & --- & --- & --- & --- \\
 & Majority     & --- & --- & --- & --- & --- \\
\bottomrule
\end{tabular}
\end{table}

Per-station results in Experiment~B (Table~\ref{tab:multi_cls_b_station}) reveal
that performance varies substantially by station. At H=24, MaxpoolCNN achieves
0.919 on CB25 and 0.910 on CB06, but only 0.827 on FD02. FD02's low macro~F1
(0.454) despite reasonable accuracy (0.827) reflects its narrow class distribution:
with only 6~active classes and two classes (2 and~3) covering 80\% of samples, the
model achieves high overall accuracy by favoring the dominant classes while
neglecting the tails.

\begin{table}[h]
\centering
\caption{Experiment~B per-station breakdown (SWC\_20 only, MaxpoolCNN).}
\label{tab:multi_cls_b_station}
\small
\begin{tabular}{llccc}
\toprule
Horizon & Station & Accuracy & Macro F1 & W-1 Acc \\
\midrule
\multirow{4}{*}{H=24}
 & FD11 & 0.882 & 0.683 & 0.981 \\
 & FD02 & 0.827 & 0.454 & 0.992 \\
 & CB06 & 0.910 & 0.674 & 0.984 \\
 & CB25 & 0.919 & 0.809 & 0.984 \\
\midrule
\multirow{4}{*}{H=168}
 & FD11 & 0.609 & 0.431 & 0.887 \\
 & FD02 & 0.675 & 0.292 & 0.968 \\
 & CB06 & 0.619 & 0.426 & 0.899 \\
 & CB25 & --- & --- & --- \\
\bottomrule
\end{tabular}
\end{table}

A notable finding from Experiment~B is SmallCNN's spatial collapse when air
temperature is added: at FD11, SmallCNN with SWC\_20+Tair drops to 0.566
accuracy at H=24 (vs.\ 0.826 with SWC\_20 only) and 0.424 at H=168 (vs.\ 0.532).
This does not occur for LargerCNN or MaxpoolCNN, suggesting that the smallest
architecture lacks the capacity to disentangle station-specific temperature
relationships during spatial generalization.

Figure~\ref{fig:multi_cls_single_vs_multi} compares single-station and
multi-station training accuracy for Experiment~A.
Figure~\ref{fig:multi_cls_holdout_station} shows the per-station accuracy
breakdown for Experiment~B.
Figure~\ref{fig:multi_cls_temporal_vs_spatial} contrasts temporal and spatial
generalization performance across all three training regimes.

\begin{figure}[h]
    \centering
    \includegraphics[width=\textwidth]{multi_cls_single_vs_multi.pdf}
    \caption{Single-station vs.\ multi-station classification accuracy
    (Experiment~A, SWC\_20 only).}
    \label{fig:multi_cls_single_vs_multi}
\end{figure}

\begin{figure}[h]
    \centering
    \includegraphics[width=\textwidth]{multi_cls_holdout_station.pdf}
    \caption{Hold-out station classification accuracy by station
    (Experiment~B, SWC\_20 only).}
    \label{fig:multi_cls_holdout_station}
\end{figure}

\begin{figure}[h]
    \centering
    \includegraphics[width=\textwidth]{multi_cls_temporal_vs_spatial.pdf}
    \caption{Temporal vs.\ spatial generalization comparison across
    single-station, Experiment~A, and Experiment~B.}
    \label{fig:multi_cls_temporal_vs_spatial}
\end{figure}

% ============================================================================
\section{Discussion}
% ============================================================================

\subsection{CNN Value is Horizon-Dependent}

The most consistent finding across all experiments is that CNNs provide meaningful
improvement over persistence specifically at long horizons (168h / 1~week ahead),
while showing no clear advantage at H=24. This is not surprising: soil moisture
changes slowly enough that the most recent observation is a strong predictor for
24~hours ahead, leaving little room for a learned model to improve. At 168~hours, the
correlation between current and future moisture has decayed substantially, and the
CNN's ability to extract patterns from the input window becomes more valuable.

\subsection{Feature Engineering Findings}

Air temperature was the only input feature that consistently improved CNN validation
performance, and only at horizons of 48~hours or longer. Seasonal sine/cosine
features, other soil moisture depths, and rainfall accumulations did not improve
validation loss. The negative result for seasonal features is worth noting: it
suggests that the day-of-year signal is already implicitly represented in the soil
moisture and temperature values themselves, so encoding it explicitly adds redundant
information rather than new predictive signal.

The failure of the multi-feature combined set (SWC at other depths + Tair) at H=24
was diagnosed as a distribution-shift issue: SWC\_5 and SWC\_10 have different value
distributions across the train/val/test splits, which destabilizes training when these
features carry substantial weight in the input.

\subsection{Architecture Observations}

Max pooling consistently outperformed average pooling, suggesting that the maximum
activation across the input window is a more informative summary than the mean
activation for soil moisture forecasting. This makes intuitive sense: extreme events
within the window (e.g., a sharp moisture increase from rainfall) may be more
predictive of future moisture state than the average pattern strength.

Larger models (more filters, more layers) did not improve test MSE, and dropout
was actively harmful. These results are consistent with a regime where the model is
already well-sized for the complexity of the patterns in the data, and additional
capacity or regularization provides no benefit.

\subsection{Generalization Across Stations and Years}

Multi-station experiments reveal that temporal generalization (predicting an unseen
year when all stations are in training) is easier than spatial generalization
(predicting at an unseen station). This pattern holds for both regression and
classification: in classification, MaxpoolCNN achieves 0.927 accuracy in
Experiment~A (hold-out year) vs.\ 0.885 in Experiment~B (hold-out station) at
H=24. The gap is larger at H=168, where temporal and spatial conditions diverge
more strongly.

Station-level differences in class distribution are a key driver of Experiment~B
performance variation. FD02, with only 6~active classes concentrated in two bins,
achieves high accuracy but low macro~F1 because the model cannot recover meaningful
per-class performance on the rare tails. CB25, with 9~active classes and the widest
SWC range, shows the strongest spatial generalization (0.919 accuracy at H=24),
suggesting that broader class coverage in the training stations provides better
transfer.

SmallCNN's spatial collapse with air temperature at FD11 (accuracy drops from
0.826 to 0.566 at H=24) indicates that the smallest architecture cannot
disentangle station-specific temperature--moisture relationships when generalizing
spatially. This failure does not occur for the two-layer architectures, pointing
to a minimum capacity threshold for spatial transfer with auxiliary features.

\subsection{Classification Observations}

The classification experiments reinforce two findings from the regression side.
First, the horizon-dependent value of learned models: at H=24, persistence
achieves 88.8\% accuracy by simply predicting that the moisture category
will not change, leaving little room for CNNs to improve (MaxpoolCNN: 89.0\%).
At H=168, persistence drops to 58.4\% while the best CNN (MaxpoolCNN) maintains
65.1\%, demonstrating clear added value at longer horizons. This pattern
strengthens with multi-station training: MaxpoolCNN reaches 0.927 accuracy at
H=24 and 0.708 at H=168 in Experiment~A, widening the gap over persistence
at the longer horizon.

Second, max pooling again outperforms average pooling across all test years,
both horizons, and both multi-station experiment types. In single-station
classification, MaxpoolCNN achieves 0.890 accuracy vs.\ LargerCNN's 0.871
at H=24. In multi-station Experiment~A, the advantage grows: 0.927 vs.\ 0.898.
In Experiment~B (spatial generalization), MaxpoolCNN averages 0.885 vs.\
LargerCNN's 0.855, confirming that max pooling's ability to detect peak
activations within the input window is especially valuable for generalization.

The distance-aware metrics reveal that classification errors are overwhelmingly
local: within-1-category accuracy exceeds 97.9\% at H=24 and 89.5\% at H=168
in single-station experiments, and remains above 98.5\% at H=24 in multi-station
Experiment~A. When the model is wrong, it is almost always wrong by only one bin
(0.05~SWC units), a practically useful property.

The class imbalance varies substantially across stations. FD11 has 9~active
classes with classes~4--8 holding $\approx$92\% of samples; FD02 has only
6~active classes with classes~2 and~3 covering 80\%. This cross-station
heterogeneity means that macro~F1 and weighted~F1 can diverge sharply:
FD02 shows 0.827 accuracy but only 0.454 macro~F1 under MaxpoolCNN at H=24,
because the two dominant classes inflate accuracy while rare-class performance
remains poor. Both metrics should be reported together to give a complete picture.

\subsection{Overfitting Analysis}

Extended training runs (100~epochs, no early stopping) show mild overfitting
emerging primarily at H=168, more pronounced for LargerCNN (+10.1\% validation loss
rise from best epoch to epoch 100) than SmallCNN (+4.7\%). LargerCNN's best
validation epoch comes consistently earlier than SmallCNN's, confirming that early
stopping is necessary for the larger model to realize its capacity advantage without
overfitting.

\subsection{Validation Loss Reporting}

An important methodological note: the \texttt{train\_model} function restores the
model to its best validation epoch via early stopping, but early results incorrectly
recorded the \emph{final} validation loss (at the last training epoch) rather than
the \emph{best} validation loss (at the restored epoch). This was corrected so that
both values are now recorded: \texttt{best\_val\_loss} (the loss at the epoch whose
weights are used) and \texttt{final\_val\_loss} (the loss at the last epoch before
stopping). Feature selection and architecture comparison use \texttt{best\_val\_loss}
for model selection.

\subsection{Limitations}

\begin{itemize}
    \item Single-seed results throughout the main experiment matrices. Small
    differences between nearby configurations should be read cautiously rather than
    as a confident ranking.
    \item The test-year wetness finding rests on a single wet year (2018). Whether
    wetness itself, or something specific to 2018, drives the higher error cannot
    be distinguished from this experiment alone.
    \item Multi-station experiments use only 4 of the 33 TxSON stations. These were
    selected for data quality, but may not be representative of the full network.
    \item The classification experiment uses a fixed 12-bin scheme (0.05-width
    intervals from 0.00 to 0.60) for all stations and years. This creates empty
    classes at each station (e.g., FD02 has six empty classes) and highly uneven
    class priors across stations. An adaptive binning scheme could equalize class
    frequencies but would sacrifice cross-station comparability.
\end{itemize}

% ============================================================================
\section{Future Work}
% ============================================================================

\begin{enumerate}
    \item Extend multi-station experiments to additional TxSON stations with
    sufficient data quality, to test whether the findings from four stations
    generalize more broadly.
    \item Implement an LSTM baseline using the same preprocessing, windowing, and
    evaluation framework, to provide a direct comparison between convolutional and
    recurrent approaches.
    \item Test temporal attention mechanisms that could preserve positional
    information lost by the current pooling step.
    \item Conduct multi-seed experiments for the strongest configurations to
    distinguish genuine performance differences from stochastic variation.
    \item Investigate alternative class-reweighting strategies for the
    classification task, such as capped inverse-frequency weights,
    $\sqrt{1/n_c}$ scaling, or focal loss, which may handle extreme class
    imbalance more gracefully than standard inverse-frequency weighting.
\end{enumerate}

\bibliographystyle{plain}
\bibliography{references}

\end{document}
